% Источники: exam/materials/преподаватель/2020/14-04-2020-Model-L9-1-14.04.20.pdf; exam/materials/преподаватель/2020/14-04-2020-Model-L9-2-14.04.20.pdf; exam/materials/моделирование.pdf, раздел 18.
\section{Имитационные модели систем массового обслуживания.}

\textbf{СМО} (система массового обслуживания) — сложная динамическая система, функционирование которой состоит в поступлении, ожидании и обслуживании заявок. Имитационная модель СМО применяется, когда аналитические формулы требуют слишком сильных предположений о потоках, дисциплине очереди или структуре обслуживания.

\subsection*{Основные понятия}

\begin{itemize}
  \item \textbf{Заявка} (транзакт) — объект, требующий обслуживания.
  \item \textbf{Источник заявок} — внешний или внутренний механизм формирования входного потока.
  \item \textbf{Обслуживающий аппарат (ОА)} — канал/прибор обслуживания.
  \item \textbf{Накопитель (очередь)} — буфер для ожидающих заявок.
  \item \textbf{Дисциплина очереди} — правило выбора следующей заявки: FIFO, LIFO, по приоритету, случайный выбор и т.п.
  \item \textbf{СМО с ожиданием} (система с накопителем) — заявки ждут в очереди.
  \item \textbf{СМО с отказами} (система с потерями) — заявка, заставшая все каналы занятыми, получает отказ и покидает систему.
  \item \textbf{Открытая СМО} — источник заявок находится вне системы.
  \item \textbf{Замкнутая СМО} — заявка после обслуживания возвращается в поток.
\end{itemize}

\subsection*{Имитационное моделирование СМО}

Имитационная модель задаёт логику движения заявок через систему и имитирует события: поступление заявки, постановку в очередь, начало обслуживания, окончание обслуживания, отказ или выход из системы. Обычно используется \textbf{дискретно-событийный} или \textbf{транзактный} подход: состояние системы меняется только в моменты событий, а статистика накапливается по траектории моделирования.

\textbf{Исходные данные модели:}
\begin{itemize}
  \item закон или эмпирическое распределение интервалов между поступлениями заявок;
  \item закон времени обслуживания на каждом аппарате;
  \item число каналов, ёмкость накопителей, маршруты заявок;
  \item дисциплина очереди, приоритеты и правила отказа;
  \item время моделирования, период разгона, число прогонов.
\end{itemize}

\textbf{Основные выходные характеристики:}
\begin{itemize}
  \item средняя длина очереди и среднее число заявок в системе;
  \item среднее и максимальное время ожидания, среднее время пребывания в системе;
  \item коэффициенты загрузки обслуживающих аппаратов;
  \item пропускная способность, вероятность отказа или потери заявки.
\end{itemize}

\subsection*{Транзактная схема GPSS}

Основные блоки GPSS:
\begin{itemize}
  \item \texttt{GENERATE} — порождает транзакты с заданным межприбытийным временем;
  \item \texttt{SEIZE} / \texttt{RELEASE} — захват/освобождение одноканального устройства (Facility);
  \item \texttt{ENTER} / \texttt{LEAVE} — захват/освобождение многоканального устройства (Storage);
  \item \texttt{QUEUE} / \texttt{DEPART} — регистрация входа/выхода из очереди для сбора статистики;
  \item \texttt{ADVANCE} — задержка транзакта на время обслуживания;
  \item \texttt{TERMINATE} — уничтожение транзакта.
\end{itemize}

\textbf{Типовая логика одноканальной СМО:}
\[
  \texttt{GENERATE} \to \texttt{QUEUE} \to \texttt{SEIZE} \to
  \texttt{DEPART} \to \texttt{ADVANCE} \to \texttt{RELEASE} \to \texttt{TERMINATE}.
\]

После серии прогонов оценивают математические ожидания, дисперсии и доверительные интервалы характеристик. Для стохастической СМО один прогон не является результатом эксперимента: нужны повторные опыты с разными последовательностями случайных чисел.

\clearpage
